% Propuesta de proyecto. Máximo 4 páginas, sin contar las referencias.
% No cambiar los títulos ni el orden de las secciones.
% Compilar con: pdflatex propuesta.tex (dos veces)
\documentclass[11pt]{article}
\usepackage[utf8]{inputenc}
\usepackage[T1]{fontenc}
\usepackage[spanish,es-noindentfirst,es-nodecimaldot,es-tabla]{babel}
\usepackage[letterpaper,margin=2.2cm,top=1.6cm]{geometry}
\usepackage{xcolor}
\usepackage{booktabs,tabularx,array,titling,enumitem,float}
\usepackage[font=small,labelfont=bf,skip=4pt]{caption}
\usepackage[hidelinks]{hyperref}

\definecolor{unamAzul}{HTML}{003D79}
\definecolor{unamOro}{HTML}{D59F0F}
\newcolumntype{L}[1]{>{\raggedright\arraybackslash}p{#1}}
\setlist{nosep,leftmargin=*}
\setlength{\parskip}{4pt}
\setlength{\parindent}{0pt}

% ---------- Encabezado institucional ----------
\pretitle{%
  \begin{center}\small
    Universidad Nacional Autónoma de México\\
    Posgrado en Ciencia e Ingeniería de la Computación\\[2pt]
    Inteligencia Artificial, semestre 2027-1
  \end{center}
  \vspace{-4pt}{\color{unamOro}\rule{\textwidth}{1.2pt}}\par\vspace{14pt}
  \begin{center}{\small\scshape Propuesta de proyecto}\\[6pt]\LARGE\bfseries\color{unamAzul}}
\posttitle{\par\end{center}\vspace{-4pt}}
\preauthor{\begin{center}\large\begin{tabular}[t]{c}}
\newcommand{\otrorenglon}{\end{tabular}\par\vspace{8pt}\begin{tabular}[t]{c}}
\postauthor{\end{tabular}\par\end{center}}
\predate{\begin{center}\small}
\postdate{\par\end{center}}

\title{GoIA: Navegando la ciudad de forma segura}
\author{%
  Belmont Hernández Héctor Antonio\\ {\small\texttt{habelmont2@gmail.com}}
  \and
  Hernández González Luis Ángel\\ {\small\texttt{luiss.angel92@gmail.com}}
  \otrorenglon
  Núñez Laces Yeni Yuzelmi\\ {\small\texttt{yenilaces@gmail.com}}
  \and
  Villa Aguilar Gerardo Erick\\ {\small\texttt{gerardoerickv@gmail.com}}
}
\date{3 de noviembre de 2026}

\begin{document}
\maketitle
\thispagestyle{empty}

% ---------------------------------------------------------------------------
% EL PROBLEMA
% ---------------------------------------------------------------------------
\section{El problema}
\label{sec:problema}

Cuando una persona es nueva en una ciudad, ya sea por mudanza o visita, una complicación común que
tienen es no saber qué zonas no son seguras para evitarlas cuando se desplazan a pie de un punto A a
un punto B. Una zona no segura se puede definir como aquella con antecedentes de crímenes como
asaltos o estafas, o condiciones como falta de iluminación que generen un estado de alerta e
incrementen las probabilidades de dichas situaciones. Una mala elección de ruta puede resultar en un
robo de cartera y celular, pudiendo llegar a un vaciado de cuentas bancarias. En otras palabras, una
de las peores experiencias para cualquier persona, con pérdida inmediata de dinero, tiempo para
reponer bienes y tarjetas, y sobre todo bienestar físico y psicológico.

Hay ciertas zonas que son ``famosas'' como Tepito y es posible obtener mucha información al respecto
en internet, pero para otros lugares más pequeños no suele haber tanta información pública y
solamente se puede saber al visitarlos en persona o hablar con los locales, cosas que requieren
llegar al lugar primero. Aplicaciones como Google Maps calculan rutas con base en el tiempo, por lo
que la ruta más rápida no es necesariamente la más segura. Hay un par de aplicaciones móviles que
serán nombradas más adelante que atacan el problema de calcular rutas considerando la seguridad pero
fueron desarrolladas fuera del país, por lo que no necesariamente se adecuan a las condiciones de
México o Latinoamérica.

La principal inspiración de nuestro proyecto es Waze, el cual le ofrece al usuario varias opciones de
rutas a escoger, y además cuenta con un sistema donde los usuarios pueden contribuir mediante
reportes. Waze no está diseñado para peatones~\cite{waze}, lo cual es un área de oportunidad.
Streetwise~\cite{streetwise} implementa un sistema similar a nuestra propuesta: un planificador de
rutas que considera la seguridad de las zonas, pero solamente está disponible en París, Lyon y
Bruselas. Además del sistema de planificación y reportes, también cuenta con un sistema de
inteligencia para patrones en los datos de la ciudad, opción que está destinada a los gobiernos.
Por otro lado, RouteSafe~\cite{routesafe} fue desarrollado en Turquía y además ofrece un asistente de
voz y un sistema de notificaciones para hacer seguimiento de la ruta de personas conocidas.

Si el problema no es solucionado, las personas seguirán con la incertidumbre de si la ruta que
eligieron es una buena opción. Solamente basta un incidente para que la vida de alguien cambie por
completo para mal. En contraste, el éxito sería una plataforma donde cualquier persona pueda tener
acceso a la información de seguridad y así tomar una decisión más informada. Este sería medido en
número de visitas a la plataforma, rutas completadas y reportes enviados.

Para el usuario cotidiano, el sistema de monetización puede ser principalmente pasivo mediante
anuncios al igual que Waze. Esto contribuiría a la idea de que la plataforma está disponible para
todos, ya que las funcionalidades principales no estarían detrás de un paywall. Dichos anuncios se
basarían en marketing para negocios, donde estos pagan para aparecer destacados con pines, logotipos
o banners patrocinados en la pantalla cuando un conductor pasa cerca de sus instalaciones.

También puede utilizarse un sistema de suscripciones que desbloquee funcionalidades de
personalización como avatares o funcionalidades de nicho como un asistente de voz. Esta sería una
idea secundaria, pues su alcance sería menor.

También sería posible colaborar con los gobiernos de las ciudades como hace Streetwise mediante
licitaciones. La propuesta de valor sería una plataforma digital moderna que el gobierno no necesita
mantener por sí mismo y con acceso a inteligencia que puede ser obtenida de los datos enviados por
los usuarios.

% ---------------------------------------------------------------------------
% SELECCIÓN DE LA IDEA
% ---------------------------------------------------------------------------
\section{Selección de la idea}
\label{sec:seleccion}

La Tabla~\ref{tab:seleccion} compara la idea elegida con las dos mejores alternativas. La idea surgió
de la operación combinar de SCAMPER, al unir las ideas de un planificador de rutas y un sistema de
seguridad ciudadana. Otras ideas que se evaluaron fueron la creación de un sistema de reportes para
deficiencias en calles como la iluminación, que son factores que influyen en la seguridad. Y si bien
es algo construible, la aplicación solo daría visibilidad al dolor de las deficiencias públicas y no
daría una solución puesto que ello depende del gobierno. Otra idea fue mostrar un mapa que muestre
zonas inseguras pero esto no necesariamente requiere IA y sería algo con bajos niveles de
monetización. El planificador de rutas seguras soluciona un dolor, hace uso de IA y es posible
obtener datos de pequeña escala para el MVP.

\begin{table}[H]
  \caption{Puntaje de 1 (bajo) a 5 (alto) de la idea elegida y de las dos mejores
    alternativas en cuatro criterios: dolor del cliente, datos disponibles, IA necesaria y
    factibilidad en 6 semanas. El total suma los cuatro criterios, y un total mayor es mejor.}
  \label{tab:seleccion}
  \begin{tabularx}{\textwidth}{@{}Xccccc@{}}
    \toprule
    Idea & Dolor & Datos & IA necesaria & Factible & Total \\
    \midrule
    Planificador de rutas seguras & 5 & 3 & 4 & 3 & 15 \\
    Sistema de reportes para deficiencias en calles & 2 & 3 & 3 & 3 & 11 \\
    Mapa que muestre zonas inseguras & 4 & 4 & 1 & 4 & 13 \\
    \bottomrule
  \end{tabularx}
\end{table}

% ---------------------------------------------------------------------------
% AGENTE Y ENTORNO DE LA TAREA
% ---------------------------------------------------------------------------
\section{Agente y entorno de la tarea}
\label{sec:agente}

El agente es el recomendador de ruta segura, que dada una ubicación permite visualizar la ruta que
tenga menor riesgo. La Tabla~\ref{tab:reas} describe la tarea con REAS, y la
Tabla~\ref{tab:propiedades} clasifica su entorno, según la caracterización de Russell y
Norvig~\cite{russell2021}.

\begin{table}[H]
  \caption{Descripción REAS de la tarea: medidas de rendimiento, entorno, actuadores y
    sensores del agente.}
  \label{tab:reas}
  \renewcommand{\arraystretch}{1.25}
  \begin{tabularx}{\textwidth}{@{}>{\bfseries}p{0.15\textwidth}X@{}}
    \toprule
    Rendimiento & Ruta seleccionada con riesgo bajo, distancia y tiempo razonable hacia el destino
                  y consideración de reportes actuales. \\
    Entorno & Calles, personas que soliciten recomendaciones, zonas urbanas, incidentes, tiempo de
              incidentes y solicitudes. \\
    Actuadores & Mostrar sobre un mapa la recomendación de una ruta con el nivel de riesgo y la
                 opción de comparar con la ruta más corta, y permitir que los usuarios ingresen
                 nuevos reportes. \\
    Sensores & Los reportes de los usuarios con coordenadas geográficas, fecha y hora, registro de
               base de datos y servicios de cartografía. \\
    \bottomrule
  \end{tabularx}
\end{table}

\begin{table}[H]
  \caption{Propiedades del entorno de la tarea, con el valor elegido para el proyecto y su
    justificación.}
  \label{tab:propiedades}
  \small
  \renewcommand{\arraystretch}{1.25}
  \begin{tabularx}{\textwidth}{@{}p{0.3\textwidth}L{0.15\textwidth}X@{}}
    \toprule
    Propiedad & Valor & Justificación \\
    \midrule
    Totalmente o parcialmente observable & Parcialmente observable & El agente conoce la
      información que ha sido ingresada por el usuario pero no conoce todos los incidentes que
      ocurren en la zona ni condiciones reales de seguridad. \\
    Un agente o multiagente & Un agente & Un solo recomendador procesa la información, evalúa los
      indicadores de riesgo y calcula la ruta segura. \\
    Determinista o estocástico & Estocástico & La ruta depende de los incidentes y condiciones de
      seguridad que no se pueden predecir con seguridad. \\
    Episódico o secuencial & Secuencial & La información recabada modifica las
      recomendaciones. \\
    Estático, semidinámico o dinámico & Dinámico & La información de seguridad va cambiando y las
      recomendaciones se van actualizando. \\
    Discreto o continuo & Discreto & El agente realiza la búsqueda sobre un conjunto finito de
      nodos restringidos por la ubicación solicitada. \\
    Conocido o desconocido & Conocido & Las reglas de búsqueda, las distancias e indicadores de
      riesgos son conocidos, aunque no se tenga información de los incidentes. \\
    \bottomrule
  \end{tabularx}
\end{table}

% ---------------------------------------------------------------------------
% HIPÓTESIS Y EVALUACIÓN
% ---------------------------------------------------------------------------
\section{Hipótesis y evaluación}
\label{sec:hipotesis}

La búsqueda de ruta mediante el algoritmo A* obtiene rutas con menor riesgo que las búsquedas
tradicionales, las cuales priorizan únicamente la minimización de tiempo, a través de indicadores de
riesgo asociados a reportes de incidentes. La Tabla~\ref{tab:hipotesis} detalla cada elemento de la
hipótesis.

\begin{table}[H]
  \caption{Elementos de la hipótesis: método, métrica, línea base y datos.}
  \label{tab:hipotesis}
  \renewcommand{\arraystretch}{1.3}
  \begin{tabularx}{\textwidth}{@{}>{\bfseries}p{0.15\textwidth}X@{}}
    \toprule
    Método & Algoritmo A* con indicadores de riesgo, tiempo y distancia de calles y zonas. \\
    Métrica & Riesgo acumulado de la ruta, calculado por los indicadores de la zona más el tiempo
              de la ruta. Se evalúa el tiempo total de traslado con una flexibilidad de desviación
              (porcentaje de tiempo adicional) a una ruta alterna. La recomendación está dada por el
              menor riesgo acumulado dado que no se sobrepase el tiempo de llegada. \\
    Línea base & Algoritmo Dijkstra o A* para obtener la ruta más rápida, sin indicadores de
                 seguridad. \\
    Datos & Utilizaremos datos sintéticos simulando zonas en las que se ha presentado un siniestro.
            Al reportar la zona se actualiza la heurística para ajustar la ruta y mostrar la más
            segura. \\
    \bottomrule
  \end{tabularx}
\end{table}

El equipo realizará simulaciones controladas para asegurarse de que el sistema responde
correctamente y evita rutas peligrosas manteniendo la alternativa más corta. Para la prueba
cualitativa, realizaremos entrevistas para validar si los usuarios preferirían las rutas dadas por
nuestro sistema.

% ---------------------------------------------------------------------------
% PRODUCTO MÍNIMO VIABLE
% ---------------------------------------------------------------------------
\section{Producto mínimo viable}
\label{sec:mvp}

El producto final sería una aplicación móvil que brinda al usuario una ruta priorizando la
seguridad y no solo la ruta más corta. Para el MVP se realizará una versión web. La
Tabla~\ref{tab:mvp} describe el producto.

\begin{table}[H]
  \caption{Descripción del producto mínimo viable: tipo, usuario, entrada, proceso y salida.}
  \label{tab:mvp}
  \renewcommand{\arraystretch}{1.3}
  \begin{tabularx}{\textwidth}{@{}>{\bfseries}p{0.15\textwidth}X@{}}
    \toprule
    Tipo & Aplicación web. \\
    Usuario & Cualquier persona que va de un origen a un destino. \\
    Entrada & El usuario proporciona su ubicación de origen y destino. El usuario ingresa reporte
              de incidentes. \\
    Proceso & El agente consulta la información de la zona con los reportes almacenados, obtiene
              indicadores de riesgo y ejecuta una búsqueda con heurística para encontrar una ruta
              combinando información de tiempo y seguridad. Si la persona tiene mucha prisa puede
              ver la ruta más rápida, compararla con la más segura y decidir cuál tomar de acuerdo
              con sus necesidades. \\
    Salida & Se muestra al usuario la ruta recomendada y sus características. \\
    \bottomrule
  \end{tabularx}
\end{table}

% ---------------------------------------------------------------------------
% PLAN DE TRABAJO
% ---------------------------------------------------------------------------
\section{Plan de trabajo}
\label{sec:plan}

Duración: 4 semanas.

El proyecto se desarrollará durante cuatro semanas, distribuyendo las responsabilidades entre los
cuatro integrantes: el primero se encargará del algoritmo A* y el cálculo de rutas; el segundo
desarrollará el mapa y la interfaz web; el tercero implementará el registro y almacenamiento de
reportes ciudadanos de inseguridad; y el cuarto integrará los componentes y realizará las pruebas.
Al finalizar la cuarta semana, pretendemos contar con un MVP funcional que permita seleccionar un
origen y destino, visualizar rutas alternativas y calcular una ruta priorizando la seguridad
mediante los reportes registrados por los usuarios.

\begin{enumerate}[leftmargin=1.5em,itemsep=4pt]
  \item Análisis de requerimientos y diseño de la solución (Semana 1, los cuatro integrantes).
    Definir el funcionamiento del MVP, los datos necesarios para representar las calles como un
    grafo, los criterios para clasificar las zonas de riesgo y el diseño preliminar de la interfaz.
    Producto: documento de requerimientos, diagrama de arquitectura y diseño inicial de pantallas.
  \item Implementación del algoritmo A* (Semanas 1 y 2, Luis). Construir la representación del mapa
    mediante nodos y aristas, implementar el algoritmo A* y establecer una función de costo que
    considere la distancia y el nivel de riesgo de cada tramo. Producto: módulo funcional de
    búsqueda de rutas seguras con datos de prueba.
  \item Desarrollo de la interfaz web y visualización del mapa (Semanas 1 y 2, Héctor). Construir la
    interfaz para seleccionar un punto de origen y destino, visualizar el mapa y mostrar las rutas
    calculadas. Producto: prototipo web navegable con mapa interactivo.
  \item Desarrollo del sistema de reportes ciudadanos (Semanas 1 y 2, Yeni). Implementar un
    formulario para registrar incidentes, incluyendo ubicación, tipo de incidente y fecha.
    Almacenar los reportes y asociarlos con los segmentos del mapa correspondientes. Producto:
    módulo de registro y consulta de reportes de inseguridad.
  \item Integración del sistema y actualización de niveles de riesgo (Semana 3, los cuatro
    integrantes). Conectar el módulo de reportes con el algoritmo A*, de manera que los incidentes
    registrados actualicen los pesos de riesgo de las calles y se reflejen en el cálculo de nuevas
    rutas. Producto: sistema integrado capaz de recalcular rutas cuando se registran nuevos
    reportes.
  \item Pruebas funcionales y corrección de errores (Semana 4, Gerardo). Evaluar el funcionamiento
    del algoritmo con diferentes escenarios, comprobar que los reportes modifiquen los costos de las
    rutas y verificar que la interfaz muestre correctamente los resultados. Producto: registro de
    pruebas, errores corregidos y versión estable del MVP.
  \item Presentación y documentación del MVP (Semana 4, los cuatro integrantes). Preparar una
    demostración del sistema y presentar los resultados obtenidos. Producto: MVP funcional y
    presentación final.
\end{enumerate}

% ---------------------------------------------------------------------------
% REFERENCIAS
% ---------------------------------------------------------------------------
\begin{thebibliography}{9}
  \bibitem{waze} Waze Community (2023). Footpaths/Walking Segments. \emph{Waze Discuss, Map Editing
    Guides}. \url{https://www.waze.com/discuss/t/footpaths-walking-segments/375447}.
  \bibitem{streetwise} Streetwise (2026). Streetwise for Cities. Streetwise.
    \url{https://www.street-wise.app/en/cities}.
  \bibitem{routesafe} NovaLabs Global (2026). RouteSafe: Safety GPS \& Family [Aplicación móvil].
    Google Play Store. \url{https://play.google.com/store/apps/details?id=app.routesafe&hl=en&pli=1}.
  \bibitem{russell2021} S. Russell y P. Norvig (2021). \emph{Artificial Intelligence: A Modern
    Approach}, 4.ª ed. Pearson. ISBN 978-0134610993. \url{https://aima.cs.berkeley.edu}.
\end{thebibliography}

\end{document}
